\documentclass[10pt,a4paper]{article}
\usepackage[margin=25mm]{geometry}
\usepackage{amsmath,amssymb,amsthm,booktabs}
\usepackage[T1]{fontenc}
\usepackage{lmodern}
\usepackage[colorlinks=true,linkcolor=blue,urlcolor=blue,citecolor=blue]{hyperref}
\usepackage{microtype}
\newtheorem{proposition}{Proposition}
\newtheorem{lemma}{Lemma}
\newcommand{\dd}{\,\mathrm d}
\newcommand{\Si}{\operatorname{Si}}
\setlength{\parindent}{0pt}
\setlength{\parskip}{6pt}
\title{Interval certification of the smooth-arc\\curvature discrepancy}
\author{Computational supplement to manuscript V423}
\date{23 September 2026}
\begin{document}
\maketitle

\textbf{Result.} The value already printed in V423 is correctly rounded to
24 decimal places:
\[
\boxed{D_\kappa\approx0.203783759536049443826081
       \quad\text{(correctly rounded to 24 decimal places).}}
\]
This is a correctly rounded approximation, not an exact finite decimal representation. The following rigorous enclosure is deliberately widened from
the full computational output:
\begin{equation}\label{eq:enclosure}
\begin{split}
0.203783759536049443826080629422606828
&<D_\kappa\\
&<0.203783759536049443826080629422606829.
\end{split}
\end{equation}
Both bounds lie strictly inside the rounding cell
\[
\bigl(
0.2037837595360494438260805,\,
0.2037837595360494438260815
\bigr).
\]
Thus the decimal retained in V423 is certified
without changing any of its 24 digits. This certificate concerns the smooth-arc
diagnostic only; it does not include the atom of turning at the Eulerian corner.

\section{The quantity and its normalisation}
Use the continuous extension at zero in
\[
r_E(\theta)=\frac{2\sin(\theta/2)}{\theta},\qquad
r_L(\theta)=b+a\cos\theta,
\]
where
\[
b=\frac{2}{\pi}\Si(\pi/2),\qquad
a=\frac{2}{\pi}\bigl(\Si(3\pi/2)-\Si(\pi/2)\bigr).
\]
Primes below denote differentiation with respect to \(\theta\). Define
\[
v_i(\theta)=\sqrt{r_i(\theta)^2+r_i'(\theta)^2},\qquad
k_i(\theta)=
\frac{r_i^2+2(r_i')^2-r_i r_i''}
     {(r_i^2+(r_i')^2)^{3/2}},
\quad i\in\{E,L\}.
\]
The signed numerator is positive on each smooth arc, so this formula equals the
unsigned curvature used in the manuscript. For the Eulerian curve, it equals
\[
\frac{3\theta\sin^2(\theta/2)+2(\theta-\sin\theta)}{\theta^3}>0
\quad(0<\theta\leq\pi),
\]
with limiting value \(13/12\) at zero. For the lima\c{c}on, its minimum is
\((b-a)(b-2a)>0\); the program verifies \(b>2a>0\).

For definiteness, both full-arc normalised arclength parameters start at
\(\theta=-\pi\) and follow increasing \(\theta\) to \(\pi\).
The profiles and speeds are even. On the positive half-arc put
\[
x=\theta/\pi,\qquad H_i(x)=\pi v_i(\pi x),\qquad
A_i(x)=\int_0^x H_i(t)\dd t,\qquad \ell_i=A_i(1).
\]
Let \(x_i(v)\) be the unique solution of \(A_i(x)=v\ell_i\), for
\(0\leq v\leq1\). Symmetry of the original full-arc integral gives exactly
\begin{equation}\label{eq:reduced}
D_\kappa^2=\int_0^1
\bigl[k_E(\pi x_E(v))-k_L(\pi x_L(v))\bigr]^2\dd v.
\end{equation}
There is no omitted factor of two: the two half-arcs contribute equally and
\(v=2u-1\) on the second half of the original parameter interval.

\begin{lemma}[Uniform inverse bound]\label{lem:inverse}
For \(0\leq x\leq1\), one has \(H_E(x)\geq2\) and \(H_L(x)>2\).
Consequently, for every exact trial point \(g\in[0,1]\),
\begin{equation}\label{eq:residual}
|x_i(v)-g|\leq\frac{|A_i(g)-v\ell_i|}{2},
\qquad
|x_i(v)-x_i(c)|\leq\frac{\ell_i}{2}|v-c|.
\end{equation}
\end{lemma}
\begin{proof}
Concavity of sine on \([0,\pi/2]\) gives \(r_E\geq2/\pi\), so
\(H_E\geq\pi r_E\geq2\).
For the lima\c{c}on, \(r_L\geq b-a>0\). Outward-rounded evaluation verifies
\(\pi(b-a)>2\), whence \(H_L\geq\pi(b-a)>2\).
The functions \(A_i\) therefore increase strictly from \(0\) to \(\ell_i\).
The first inequality follows from the mean value theorem; the second follows
from \(x_i'(v)=\ell_i/H_i(x_i(v))\).
\end{proof}

Newton iteration is used only to obtain an accurate trial point.
The interval residual and Lemma~\ref{lem:inverse}, rather than the iteration's
stopping rule, establish the inverse enclosure.

\section{Enclosing the removable singularity and all derivatives}
The numerical computation avoids cancellation in the sinc formula at zero.
It uses the entire series
\[
r_E(\theta)=\sum_{k=0}^{\infty}
  \frac{(-1)^k\theta^{2k}}{2^{2k}(2k+1)!}.
\]
Let \(R_m\) retain \(k=0,\ldots,m\), and set \(n=2m+2\).
For \(|\theta|\leq4\) and \(0\leq j<n\), define
\[
q_j=\frac{4}{(n+2-j)(n+1-j)}.
\]
When \(q_j<1\), absolute summation of the differentiated tail gives
\begin{equation}\label{eq:sinc-tail}
\frac{|(r_E-R_m)^{(j)}(\theta)|}{j!}
\leq
\frac{4^{\,n-j}}
 {2^n(n+1)(n-j)!\,j!}\,\frac{1}{1-q_j}.
\end{equation}
Indeed, the ratio of successive absolute differentiated terms is at most
\[
\frac{\theta^2}{4}
\frac{n+1}{n+3}
\frac{1}{(n+2-j)(n+1-j)}
\leq q_j,
\]
and this upper bound decreases along the tail.

The implementation uses \(m=80\). It shifts the finite polynomial to the real
point or interval under consideration and adds the bound
\eqref{eq:sinc-tail} to each derivative coefficient.
The coefficients for \(r_E'\) and \(r_E''\) follow by differentiation, including
the corresponding tail bounds. All evaluation centres are checked to satisfy
\(|\theta|<4\).

Real ball arithmetic then encloses the Taylor coefficients of \(H_i\) and \(k_i\)
under addition, multiplication, division and the positive square root.
Every speed denominator is checked to have a strictly positive enclosure.
For the lima\c{c}on, sine and cosine power-series operations supply the
coefficients directly.
The series truncation symbol alone is not used as an error estimate:
\eqref{eq:sinc-tail} and the explicit quadrature remainders below provide the
required bounds.

\section{Validated arclength and outer quadrature}
Fix an even order \(p\) and a dyadic partition of \([0,1]\) into \(N\) panels.
The principal run uses \(p=32\), \(N=64\), and 256-bit midpoint precision.

\textbf{Arclength integration.}
For a panel \([a,a+h]\), interval series evaluation over the whole panel encloses
\(H_i^{(p)}(\xi)/p!\). If its absolute value is bounded by \(B\), then for
\(0\leq d\leq h\),
\begin{equation}\label{eq:inner}
A_i(a+d)-A_i(a)
\in
\sum_{j=0}^{p-1}\frac{H_i^{(j)}(a)}{j!}\frac{d^{j+1}}{j+1}
+
\left[-\frac{B d^{p+1}}{p+1},\frac{B d^{p+1}}{p+1}\right].
\end{equation}
Accumulating these enclosures gives both \(A_i(g)\) and \(\ell_i\).
Coefficient rounding errors and all previous-panel errors remain in the balls.

\textbf{Inverse derivatives.}
At a base point \(x_0\), form the series
\[
F_i(s)=\frac{A_i(x_0+s)-A_i(x_0)}{\ell_i}.
\]
Its constant term is exactly zero and its linear coefficient is positive.
Formal series reversion therefore gives the derivatives of
\(x_i(v_0+s)-x_0\), where \(v_0=A_i(x_0)/\ell_i\).
Composing with the curvature series gives derivatives of
\(k_i(\pi x_i(v))\).
For interval \(x_0\), the same ball operations enclose these derivatives for
every base point in the interval.
The residual enclosure determines \(x_0\) at the panel centre.
The second bound in \eqref{eq:residual} covers all inverse images of the panel.

\textbf{Outer integration.}
Write \(f(v)\) for the squared difference in \eqref{eq:reduced}.
On a symmetric panel \([c-h,c+h]\), let \(a_j\) enclose \(f^{(j)}(c)/j!\),
and let \(B_p\) bound \(|f^{(p)}(v)|/p!\) throughout the panel.
Taylor's theorem gives
\begin{equation}\label{eq:outer}
\int_{c-h}^{c+h}f(v)\dd v
\in
\sum_{\substack{0\leq j<p\\j\ {\rm even}}}
  a_j\,\frac{2h^{j+1}}{j+1}
+
\left[-\frac{2B_p h^{p+1}}{p+1},
       \frac{2B_p h^{p+1}}{p+1}\right].
\end{equation}
The interval inverse series on the whole panel supplies \(B_p\).
Summing \eqref{eq:outer} encloses \(D_\kappa^2\); the positive square root
encloses \(D_\kappa\).

\begin{proposition}[Certification criterion]
Under the checked positivity and domain conditions above, the interval returned
by the supplied program contains the value in \eqref{eq:reduced}.
If that interval lies strictly inside a decimal's nearest-rounding cell, the
decimal is correctly rounded.
\end{proposition}
\begin{proof}
The entire-series bound encloses the radial derivatives.
Ball arithmetic propagates those inclusions through the speed and curvature
formulae. Equation~\eqref{eq:inner} encloses every arclength evaluation,
Lemma~\ref{lem:inverse} encloses every inverse value, and series reversion
encloses its derivatives. Equation~\eqref{eq:outer} includes the complete
quadrature remainder. Summation and the monotone square root preserve inclusion.
Strict containment in a rounding cell excludes either possible rounding boundary.
\end{proof}

\section{Computed certificate and checks}
The principal run returns the following enclosure for \(D_\kappa\):
\begin{align*}
L={}&0.203783759536049443826080629422606828470494451196624581353928,\\
U={}&0.203783759536049443826080629422606828474246989273842126201317.
\end{align*}
The machine-readable record contains the complete enclosure for the square,
the two half-lengths, the numerical parameters and the exact program hash.
The wider interval \eqref{eq:enclosure} follows immediately.

\begin{center}
\begin{tabular}{lrrr}
\toprule
Run & Precision (bits) & Order \(p\) & Panels \(N\)\\
\midrule
Principal certificate & 256 & 32 & 64\\
Finer certificate & 320 & 40 & 128\\
\bottomrule
\end{tabular}
\end{center}
The finer enclosure is contained in the principal enclosure. Both certify the
same 24-place value, and both also pass the program's 36-place rounding test.
The certificate is the enclosure and its error argument, not agreement between runs.

An independent 60-decimal-place computation uses direct Gauss--Legendre
quadrature and Newton inversion of separately integrated arclength functions.
Increasing the outer quadrature order from 32 to 48 to 64 converges towards the
certified value; its order-64 estimate lies inside the principal enclosure.
This is supplementary numerical corroboration, not an interval certificate.
An exploratory ordinary ODE solver did not attain its requested accuracy and was
rejected as evidence.

The manuscript's existing decimal therefore needs no numerical change.
Its accompanying accuracy statement can now cite validated arclength inversion
and Taylor quadrature. The normalisation should explicitly identify the common
starting angle and orientation.

\section{Reproduction and scope}
The primary dependency is \texttt{python-flint==0.9.0}, which supplies FLINT/Arb
ball and power-series arithmetic~\cite{arb,series}.
The optional independent check uses \texttt{mpmath==1.4.1}.
The certificate was run with Python 3.12.14.

From the supplement directory:
\begin{verbatim}
python3 certify_dk.py --order 32 --panels 64 \
  --output D_K_CERTIFICATE_32_64.json
python3 certify_dk.py --bits 320 --order 40 --panels 128 \
  --output D_K_CERTIFICATE_40_128.json
python3 verify_certificate.py
\end{verbatim}
The first two commands fail if the returned enclosure cannot certify the
manuscript's 24-place value. The last command checks the recorded enclosures,
their rounding-cell inclusion, the independent comparison, and the program hash.

This is a rigorous numerical certificate conditional on the documented
floating-point ball-arithmetic implementation and its standard software/hardware
assumptions; it is not a proof-assistant formalisation.
It certifies the stated smooth-arc scalar diagnostic, not a global equality of
curvature measures and not the entire manuscript.
The numerical definitions and displayed values are unchanged in V423;
its accompanying accuracy statements now refer to this certificate.

\begin{thebibliography}{9}
\bibitem{arb}
F. Johansson, \emph{Arb: Efficient Arbitrary-Precision Midpoint-Radius Interval
Arithmetic}, arXiv:1611.02831.
\url{https://arxiv.org/abs/1611.02831}.
See also the current FLINT documentation:
\url{https://flintlib.org/doc/index_arb.html}.
\bibitem{series}
\emph{python-flint 0.9.0 documentation: arb and arb\_series}.
\url{https://python-flint.readthedocs.io/en/latest/arb.html};
\url{https://python-flint.readthedocs.io/en/latest/arb_series.html}.
\end{thebibliography}
\end{document}
